\clearpage
\VerdierPhysicalPageAnchor{0119}%
\begingroup
\small
\setlength{\abovedisplayskip}{4pt plus 2pt minus 2pt}%
\setlength{\belowdisplayskip}{4pt plus 2pt minus 2pt}%

\VerdierSemanticAnchor{ch2-1-2-8-splitting-complete}
$fgf=f$ et $gfg=g$. Un morphisme est dit \emph{scindable} s'il possède un
scindage. Les morphismes décomposables sont scindables. La réciproque n'est
pas nécessairement vraie (exemple : catégorie des modules libres sur un
anneau de Dedekind non principal).

\VerdierSemanticAnchor{ch2-1-2-8-decomposable-category}
Une catégorie additive est dite \emph{décomposable} si tous les morphismes
de cette catégorie sont décomposables. Une catégorie décomposable est
abélienne.

\VerdierSemanticAnchor{ch2-1-2-8-triangle-decomposition}
Soit $X\xrightarrow{f}Y$ un morphisme d'une catégorie triangulée
$\mathcal D$ et soit :

\VerdierSemanticAnchor{ch2-1-2-8-decomposition-diagram-repeat}
\begin{equation*}
  \begin{tikzcd}[ampersand replacement=\&,column sep=5.2em,row sep=3.4em]
    X \arrow[r,"f"] \arrow[d,"\wr"']
      \& Y \arrow[d,"\wr"] \\
    Z\oplus W
      \arrow[r,"{\left(\begin{smallmatrix}0&\operatorname{id}_W\\0&0\end{smallmatrix}\right)}"]
      \& W\oplus Z'
  \end{tikzcd}\,.
\end{equation*}
une décomposition de $f$. Tout triangle distingué contenant $f$ est
isomorphe au triangle :

\VerdierSemanticAnchor{ch2-1-2-8-decomposed-triangle}
\begin{equation*}
  \begin{tikzcd}[ampersand replacement=\&,column sep=2.0em]
    Z\oplus W
      \arrow[r,"{\left(\begin{smallmatrix}0&\operatorname{id}_W\\0&0\end{smallmatrix}\right)}"]
      \& W\oplus Z'
      \arrow[r,"{\left(\begin{smallmatrix}0&\operatorname{id}_{Z'}\\0&0\end{smallmatrix}\right)}"]
      \& Z'\oplus Z[1]
      \arrow[r,"{\left(\begin{smallmatrix}0&\operatorname{id}_{Z[1]}\\0&0\end{smallmatrix}\right)}"]
      \& Z[1]\oplus W[1]
  \end{tikzcd}\,.
\end{equation*}

\VerdierSemanticAnchor{ch2-1-2-8-decomposed-triangle-derivation}
Ceci résulte immédiatement de (1.2.4), de (1.2.5) et des axiomes TRI et
TRII.

\VerdierSemanticAnchor{ch2-1-2-9}
\medskip
\noindent\textbf{Proposition 1.2.9.}\quad
\emph{Soient $\mathcal D$ une catégorie triangulée et
$f\colon X\longrightarrow Y$ un morphisme de $\mathcal D$. Considérons les
propriétés suivantes du morphisme $f$ :}

\VerdierSemanticAnchor{ch2-1-2-9-properties}
\begin{enumerate}
\renewcommand{\labelenumi}{\roman{enumi})}
\item \emph{Le morphisme $f$ admet un noyau.}
\item \emph{Le morphisme $f$ admet un conoyau.}
\item \emph{Le morphisme $f$ est décomposable.}
\item \emph{Le morphisme $f$ est scindable.}
\end{enumerate}

\noindent\emph{On a alors les implications :}

\VerdierSemanticAnchor{ch2-1-2-9-implications}
\[
  \mathrm{i)}\Longleftrightarrow\mathrm{ii)}
  \Longleftrightarrow\mathrm{iii)}\Longrightarrow\mathrm{iv)}\,.
\]

\VerdierSemanticAnchor{ch2-1-2-9-countable-equivalence}
\emph{Lorsque les produits dénombrables ou bien lorsque les sommes
dénombrables sont représentables dans $\mathcal D$, les propriétés
\textnormal{i)}, \textnormal{ii)}, \textnormal{iii)} et \textnormal{iv)}
sont équivalentes.}

\VerdierSemanticAnchor{ch2-1-2-9-proof-opening}
Tout d'abord, on a évidemment les implications
$\mathrm{iii)}\Longrightarrow\mathrm{i)}$,
$\mathrm{iii)}\Longrightarrow\mathrm{ii)}$,
$\mathrm{iii)}\Longrightarrow\mathrm{iv)}$. L'implication
$\mathrm{i)}\Longrightarrow\mathrm{iii)}$ est équivalente à l'implication
\endgroup
